\section{Generalisation}\label{sec:general}

\subsection{A public GPT-2 SAE family}\label{sec:e2}

To check that our conclusions are not an artefact of small, in-house SAEs, we applied the same evaluation to the public
\texttt{gpt2-small-res-jb} SAEs (24{,}576 latents, 300M training tokens) at layers 8 and 10, with layers 1 and 11 as a
labelled extension (\evlabel{validation}; Appendix Table~\ref{tab:e2}). The pattern holds: the frozen truth loss is
mostly reader fragility, the persona loss is real and survives re-training, the SAE does not beat PCA at equal
reconstruction error, and the reconstruction carries only part of the persona effect at layers 8 and 11 (0.56 and 0.61).
The one exception is the re-trained truth gap at layer 8, which is small with an interval touching zero, as for our own
SAEs. This is a single public family, so it shows that the conclusions do not depend on our training recipe, not that
they hold for every SAE.

\subsection{Gemma-2-2B: behavioural support}\label{sec:gemma}

The second model tests whether the behavioural finding depends on GPT-2 and on our own SAEs. The experiment was
pre-registered with model, tokenizer, SAE and data revisions pinned and evaluated on the Pile validation file
(\evlabel{validation}). At each of layers 5, 12 and 19, the Gemma Scope SAE reconstructs the activations of 512
documents (254{,}144 non-BOS tokens per layer), and the reconstruction is spliced back at that layer only. The
pre-registered hypothesis H-G1, carried over from GPT-2, states that SAE splicing raises next-token cross-entropy more
than a norm-matched random-error control; it is supported at a layer if the lower bound of the 95\,\% document-bootstrap
interval of the mean paired difference is above zero. All five pre-registered validity gates passed, including an
identity splice and a repeated pass that changed the output distribution by exactly zero (Appendix~\ref{app:m6}).

At their released sparsity the SAEs explain 0.813, 0.770 and 0.789 of the variance, and our measured L0 is within
4\,\% of the \texttt{average\_l0} values published with the release (Table~\ref{tab:m6}). Splicing raises
cross-entropy by 0.071, 0.115 and 0.177 nats per prediction, from an original 1.938, whereas the random-error control
raises it by 0.029, 0.039 and 0.080 nats. The mean paired difference is 0.041, 0.077 and 0.099 nats, with 95\,\%
intervals [0.036, 0.046], [0.073, 0.082] and [0.094, 0.105] (Figure~\ref{fig:m6-hg1}). H-G1 is supported at all three
layers, also under a Bonferroni interval: SAE reconstruction error produced greater next-token damage than the
norm-matched random-error control at all three pre-registered layers. The experiment does not identify the mechanism.

This result supports the behavioural finding only. Gemma was not evaluated at the information or causal levels, so
nothing here bears on information or causal preservation in Gemma, and because the two settings differ in model, SAE
family, width, sparsity, hook, data and context length, no comparison of the models is made. Each Gemma layer has its
own SAE, so the layers are not ranked.

% written by a one-off script from results/gemma2_pile/gemma_scope_splicing/summary_by_layer.csv and h_g1.csv
% (values read from the CSVs, not typed by hand); do not edit numbers by hand.
\begin{table}[t]
\centering
\small
\caption{M6 (\evlabel{validation}): Gemma Scope width-16k JumpReLU SAEs spliced into Gemma-2-2B at \texttt{resid\_post}
on 512 Pile validation documents (254{,}144 non-BOS tokens per layer). L0: the release's own \texttt{average\_l0} (published
with the SAE) / our measurement on this sample. $\Delta$CE in nats per prediction (original CE 1.938); brackets are 95\,\%
document-bootstrap intervals (1{,}000 resamples), not seed variability. $\bar d$: mean over documents of
$\Delta$CE$_{\text{SAE}}-\Delta$CE$_{\text{random}}$ (H-G1). Layers are described side by side, not ranked; each layer has its
own SAE. Source: \texttt{results/\allowbreak{}gemma2\_pile/\allowbreak{}gemma\_scope\_splicing/}.}
\label{tab:m6}
\begin{adjustbox}{max width=\linewidth}
\begin{tabular}{rcccccccl}
\toprule
Layer & L0 (release / measured) & EV $\uparrow$ & $\Delta$CE$_{\text{SAE}}$ $\downarrow$ & $\Delta$CE$_{\text{random}}$ & KL $\downarrow$ & Top-1 agr.\ $\uparrow$ & $\bar d$ [95\,\% CI] & H-G1 \\
\midrule
5 & 68 / 69.5 & 0.813 & 0.071 [0.064, 0.079] & 0.029 [0.026, 0.032] & 0.082 & 0.887 & 0.0410 [0.0356, 0.0463] & supported \\
12 & 82 / 84.8 & 0.770 & 0.115 [0.107, 0.124] & 0.039 [0.034, 0.044] & 0.120 & 0.857 & 0.0773 [0.0727, 0.0822] & supported \\
19 & 73 / 75.1 & 0.789 & 0.177 [0.165, 0.192] & 0.080 [0.069, 0.094] & 0.179 & 0.838 & 0.0993 [0.0938, 0.1051] & supported \\
\bottomrule
\end{tabular}
\end{adjustbox}
\end{table}


\begin{figure}[t]
  \centering
  \includegraphics[width=\linewidth]{fig3_h_g1_paired_difference}
  \caption{H-G1 in Gemma-2-2B (\evlabel{validation}). Left: per-document paired difference
  $d=\Delta\mathrm{CE}_{\text{SAE}}-\Delta\mathrm{CE}_{\text{random}}$ for 512 Pile validation documents (boxes show
  quartiles). Right: mean $d$ with 95\,\% (thick) and Bonferroni 98.33\,\% (thin) document-bootstrap intervals. Each
  layer has a different SAE, so values are not compared between layers.}
  \label{fig:m6-hg1}
\end{figure}
